\chapter{Room Ballot}\label{App:RoomBallot}

\section{Room Ballot General}

\appnpara{The overall order of choice on the Room Ballot shall be as follows: the Fourth Year Ballot, then the Third Year Ballot, with the Second Year Ballot at the bottom.}
\npara A timeline for the Ballot and Ballot Exemption processes (the ‘Ballot Timeline’), based on the previous year’s timeline, should be agreed between the JCR Vice-President, Head of Accommodation and Academic Registrar by the end of Michaelmas Term.  
\npara Should a member not wish to live in College Accommodation in the following academic year, they must inform College and the Vice President by the established deadline in Hilary Term.

\section{The Second Year Ballot}

\npara The Second Year Ballot will be made up of current first year students and any suspended student who chooses to enter this ballot (as per point \ref{AppP:SuspendedStudents})
\npara Students may choose to enter the ballot individually or in groups of up to six. The Vice-President shall solicit for group entries in Hilary Term. Failure to respond to this email before the deadline will be treated as a desire to enter the ballot as an individual.
\npara The order of the Second Year Ballot will be determined by drawing the names of students out of a hat, with each entry (individual or group) written on a single piece of paper. This is to be carried out in Hilary Term in the JCR at an established date and time. The order of students within each group shall be determined by a mini-draw. The JCR Vice-President may do this in their own time, either before or after the Ballot Order is drawn.
\npara The finalised order of the Second Year Ballot will be made available to all JCR members within 24 hours of the draw.

\section{The Third Year Ballot}

\npara The Third Year Ballot will be made up of current second year students and any suspended student (as per \ref{AppP:SuspendedStudents}), fourth year Law Course II student, or fourth year Modern Linguist (as per \ref{AppP:ModLing}) who chooses to enter this ballot.

\npara \label{AppP:ModLing}Current third year Law Course II students and Modern Linguists returning from their year abroad may choose to either join the Fourth Year Ballot or the Third Year Ballot, for the latter taking up the position they would have had if they had been in residence in their third year. This choice will be offered via email by the Vice President. Failure to respond by the deadline in Hilary Term will be taken as a desire to enter the Fourth Year Ballot.

\npara The order of the Third Year Ballot will be established by the following steps strictly in this order:
\begin{enumerate}
\item The previous year’s Second Year Ballot is inverted.
\item  	Students joining the ballot at a position determined by random number generator (as per \ref{AppP:RandomNumInsert} and \ref{AppP:ExceptionalRNI}) are inserted. 
\item Law Course II students and Modern Linguists returning for their fourth year who have chosen to re-enter the Third Year Ballot are inserted at the position they would have occupied in their third year. \\ \\
For example, if Walter de Merton (a fourth year linguist, returning from his year abroad) would have been 3rd on the Third Year Ballot if he had been in residence in third year, he is assigned ‘Heads’ and the current no. 3 (Peter of Abingdon, in our example) on the Ballot is assigned tails. A coin is then flipped, with the winner taking position 3 and the loser position 4. The ballot order would look as follows if tails won: \\ \\
1. Richard Werplysdon \\
2. John de la More \\
3. Peter of Abingdon \\
4. Walter de Merton (re-inserted 4th year Modern Linguist) \\
5. Robert Trenge \\
6. William Durant \\
\item  Law Course II students and Modern Linguists who will be on their year abroad in 3rd year, students who have chosen to live in non-college accommodation, and students who have been granted Ballot Exemption (who participated in the ballot in the previous year) are removed.
\end{enumerate}

\section{The Fourth Year Ballot}

\npara The Fourth Year Ballot will be made up of current third year students who are studying a four year integrated masters or Classics degree, and any fourth year Law Course II student, fourth year Modern Linguist (as per \ref{AppP:ModLing}) or suspended student (as per \ref{AppP:SuspendedStudents}) who chooses to enter this ballot.
\npara Students may choose to enter the ballot individually or in groups of up to six. The Vice-President shall solicit for group entries in Hilary Term. Failure to respond to this email before the deadline will be treated as a desire to enter the ballot as an individual.
\npara The order of the Fourth Year Ballot will be determined by drawing the names of students out of a hat, with each entry (individual or group) written on a single piece of paper. This is to be carried out in Hilary Term in the JCR at an established date and time. The order of students within each group shall be determined by a mini-draw. The JCR Vice-President may do this in their own time, either before or after the Ballot Order is drawn.
\npara The finalised order of the Fourth Year Ballot will be made available to all JCR members within 24 hours of the draw.

\section{Suspended Students} 

\npara \label{AppP:SuspendedStudents} Suspended students will be kept in the Ballot of their matriculation year for as long as possible.

\npara \label{AppP:SuspendedAcademicYear} If it is not possible for the suspended student to participate in a Ballot with their matriculation year (ie. this year have graduated) or if the suspended student would prefer, they should be placed in the Ballot relevant to their academic year.

\npara \label{AppP:RandomNumInsert} If  \ref{AppP:SuspendedAcademicYear} would result in the student participating in the Third Year Ballot but they did not have a position in the Second Year Ballot the previous year, the student should be inserted in the order at a position determined by a random number generator. \\ 
To insert by random number generator, the process will be as follows: 
\begin{enumerate}
\item Identify the total number of students currently in the ballot (excluding any students due to be inserted into the ballot), N.
\item Choose a random number, k, between 1 and N + 1.
\item Insert the student into the kth position, moving lower students down.
\end{enumerate}
\npara The Vice-President will communicate via email with all suspended students returning in the next academic year to explain their options and to confirm their choice. Failure to respond to this email by the established deadline will give the Vice-President the power to place the student in the ballot which they deem most appropriate.

\section{Exceptional Circumstances}
\npara \label{AppP:ExceptionalRNI} In the exceptional circumstance (not covered by any of the previously established processes) whereby a student was not included in the previous year’s ballot but is to be included in this year’s ballot (for example, if a student was assigned a room through ballot exemption in the previous year, but is not being assigned a room through this system this year), the student should be inserted into the ballot of their matriculation year (if not possible, their academic year). If this results in them partaking in the Third Year Ballot, the process is the same as \ref{AppP:RandomNumInsert}

\npara Any students on a course with a year abroad in second year (for example, ab initio Russian students) will be treated as follows. They should be given a position in the Second Year Ballot as normal, as if they will be in residence the following year, but removed once the order is determined. They should be reinserted into the Third Year Ballot the following year before the order is inverted. In fourth year, they will participate in the Fourth Year Ballot.

\section{Room Choosing}

\npara The Ballot Document (a list of all rooms with reviews, photos and floor plans) should be updated by the Vice-President and then circulated to the JCR by, at the latest, the end of Hilary Term. This should be updated before the Room Visits to indicate which rooms will be available for choosing that year, which rooms have been assigned already via the Ballot Exemption process, and which houses are to be assigned ‘Fourth Year Houses’.

\npara A Ballot Spreadsheet will be released by Sunday 1st Week of Trinity Term, where students can indicate their preferences and choices before the Room Choosing calls.

\npara The Head of Accommodation will assign a number of Holywell houses to be ‘Fourth Year Houses’. Participants in the Fourth Year Ballot can only choose from these rooms. The total number of rooms in these houses shall not be less than the number of participants in the Fourth Year Ballot, however it may exceed the number of participants. In this case, the remaining rooms will be released to the Third and Second Year Ballots after the Fourth Year Ballot Room Choosing.

\npara Room Visits will take place in the 1st week of Trinity Term.

\npara Room Choosing for all Ballots should take place before the end of Trinity Term. 

\npara If a member will not be present at their Room Choosing call, they may give permission to choose their room to another member of their Ballot. They must notify the Vice-President of this before the call.

\npara If a member does not attend their Room Choosing call and has not notified the Vice-President of their nominated chooser, the Vice President shall choose what they consider to be the most suitable room on behalf of that member. 

\npara In the event of a room becoming available after room choosing, this will be advertised to the JCR and the room will be made available to the person highest on the room ballot who wishes to move into that room.

